\documentclass[9pt,twocolumn,twoside]{pnas-new}
\templatetype{pnasresearcharticle}
\usepackage{pgfplots}
\pgfplotsset{compat=1.17}
\usetikzlibrary{arrows.meta}
\usepackage{makecell}
\title{A model of optimal science funding: targeting the capability-resource gap}
\author[a,1]{Aydin Mohseni}
\author[b]{Simon DeDeo}
\author[a]{Kevin J.\,S. Zollman}
\affil[a]{Department of Philosophy, Carnegie Mellon University, Pittsburgh, PA 15213}
\affil[b]{Department of Social and Decision Sciences, Carnegie Mellon University, Pittsburgh, PA 15213}
\leadauthor{Mohseni}
\significancestatement{Research funders divide fixed budgets among researchers whose abilities and needs they cannot observe. Whether to concentrate funds or spread them, whether peer review is worth its cost, and whether lotteries should replace it are unsettled. We model funding as a sequential decision under uncertainty and show that the output-maximizing allocation funds the gap between each researcher's capability and their resources, that track records alone underdetermine that gap, and that what peer review is worth and how lotteries and seed grants perform turn on two features of a field: the funder's budget relative to researchers' resources, and how unequal the gaps are, which mostly reflects how unequal capability is. We show where representative funders fall on these two features.}
\authorcontributions{A.M., S.D., and K.J.S.Z. designed research; A.M. performed research and analyzed data; A.M., S.D., and K.J.S.Z. wrote the paper.}
\authordeclaration{The authors declare no competing interest.}
\correspondingauthor{\textsuperscript{1}To whom correspondence should be addressed. E-mail: amohseni@andrew.cmu.edu}
\keywords{science funding $|$ peer review $|$ lotteries $|$ Bayesian decision theory $|$ metascience}
\begin{abstract}
The grant funder's challenge is to maximize the impact of a fixed budget under uncertainty about the researchers who might receive it. We develop a Bayesian model of this challenge: researchers differ in their capabilities and in their resources; research output requires both, and funding compounds capability across rounds. We show that the output-maximizing allocation of a round's budget funds the capability-resource gap: grants increase with capability and decrease with existing resources. Track records underdetermine the gap, because output confounds the two; what a funder needs is a signal of capability separate from resources, and peer review is one source of such a signal. We find that the value of that signal increases with how unequally the gaps are distributed, which capability inequality mainly drives, and with the signal's informativeness, and that under heavy-tailed capability distributions even a noisy signal captures most of that value. Uniform seed grants and lotteries, which forgo targeting, lose little research output exactly where review is worth little, and most of what a partial lottery loses comes from discarding the review ranking rather than from chance. Two features of a field, the funder's budget relative to researchers' resources and the inequality of the gaps, locate a funder on a map of these results.
\end{abstract}
\dates{This manuscript was compiled on \today}
\doi{\url{www.pnas.org/cgi/doi/10.1073/pnas.XXXXXXXXXX}}
\begin{document}
\maketitle
\thispagestyle{firststyle}
\ifthenelse{\boolean{shortarticle}}{\ifthenelse{\boolean{singlecolumn}}{\abscontentformatted}{\abscontent}}{}
\dropcap{T}he United States spent close to a trillion dollars on research and development in 2024, roughly a fifth of it federal \citep{NCSES2026, NSB2026}, and how best to allocate that money remains unsettled. Whether funding should be concentrated among a few researchers or spread across many is a long-running dispute \citep{Aagaard2020}. How well peer review predicts research outcomes is contested: among funded NIH grants, percentile scores predicted subsequent productivity barely better than chance \citep{Fang2016}, while studies with other designs find modest but real predictive power \citep{LiAgha2015, ParkLeeKim2015}. Several funders now allocate part of their budgets by lottery, from a tie-break among top-ranked proposals to a draw among all applicants above a threshold \citep{Liu2020, Heyard2022, Feliciani2024}. And a randomized trial of winning a grant found no clear effect on the recipients' publications over four years \citep{Barnett2024}, which raises the prior question of when a marginal grant adds output at all.

These disputes concern one problem: what a funder with a fixed budget and imperfect knowledge of researchers should do. Each empirical study measures one input to that decision; the answer depends on what the funder is uncertain about, what its information is worth, and how the budget changes both. We treat the funder's problem as a sequential decision under uncertainty and analyze it with a Bayesian approach. Researchers differ in capability, what they know and can do, and in resources, what they have to work with. Output requires both and is limited by the scarcer of the two. The funder observes output and the noisy judgments of reviewers, but neither attribute directly. A grant adds resources in the round it is given and, because doing research builds expertise, raises capability in later rounds. Each round, the funder observes the output, updates its beliefs, and allocates again.

Three results follow. First, the allocation of a round's budget that maximizes expected research output funds the gap between each researcher's capability and their resources: grants increase with capability and decrease with existing resources, and the budget sets how many researchers are funded at all. Two heuristics a funder might use, funding the track record and funding the under-resourced, each track one of the two quantities, and each can perform worse than dividing the budget uniformly across all researchers. Second, track records underdetermine the gap, because output confounds capability with resources; what the funder needs is a signal of capability separate from resources, and peer review is one source of such a signal. The value of that signal, and how seed grants and lotteries perform relative to targeted funding, depend on two features of a field: how much of researchers' resources the funder's budget can supply, and how unequally the capability-resource gaps are distributed, which in every setting we ran is mostly a matter of how unequal capability is. Where the budget is tight and capability heavy-tailed, targeting is worth a great deal, even noisy review captures most of that value, and lotteries give up much of it; where the budget is ample and capability evenly spread, little research output is at stake and spreading funds widely loses almost nothing. Third, the two features can be estimated from what funders already know, and we place representative funding programs on the map they define.

Two assumptions bound these results: grants are used up in the round they are given, with whatever lasts from them entering as capability, and capability and resources are complements. The objective is aggregate expected research output; diversity of directions, risk, and fairness are outside it, and ``optimal'' throughout means output-maximizing under that objective. The model is deliberately simple. It shows that the disputes turn on two features of a field, and which way each dispute goes in each kind of field.

\section*{Results}

\subsection*{The optimal allocation funds the capability-resource gap}
Consider first a funder that knows every researcher's capability $K_i$ and resources $R_i$ and divides one round's budget $B$ so as to maximize that round's expected output, which for researcher $i$ is $AK_iR_i/(K_i + R_i)$, where $A$ sets the units of output and $R_i$ is the sum of baseline resources $R_{i0}$ and the grant $g_i$ (Materials and Methods). The marginal value of a dollar to researcher $i$ is $AK_i^2/(K_i + R_{i0} + g_i)^2$: it increases with their capability and decreases with the resources they already hold, grants included. An allocation is optimal exactly when no dollar can be moved to a researcher who would produce more with it, and that condition gives the gap rule (SI Text 1, Proposition~1):
\[
g_i^* = \max(c\,K_i - R_{i0},\, 0),
\]
where $c > 0$ is the value at which the grants sum to the budget. Each researcher's target is $c$ times their capability; the rule fills each positive gap between target and resources and grants nothing to the rest. The rule holds for any production function with constant returns to scale and diminishing returns to resources, $F(K, R) = A\,K\,h(R/K)$ with $h$ increasing and strictly concave, because equal marginal value among the funded then means a common ratio of resources to capability (SI Text 1, Proposition~1); the harmonic mean is the instance we simulate. The condition $R_{i0} < cK_i$ draws a frontier through the population, separating the funded from the unfunded (Fig.~\ref{fig:p1}). A larger budget raises $c$, moves the frontier outward, and funds more researchers with larger grants, until funding extends across nearly the whole population.

The gap rule explains why each heuristic fails. Funding the track record fails because output reflects resources as well as capability: a productive, well-resourced researcher has a small or negative gap, and an additional dollar adds little to their output. Funding the under-resourced fails for the opposite reason: scarcity is no evidence of capability, and dollars given to researchers of low capability add little output however scarce their resources. Each heuristic can produce less output than dividing the budget equally (SI Text 1, Examples 1 and 2; Fig. S1).

How much targeting matters depends on the budget. Expected output never exceeds $AK_i$, so when the budget is large enough to bring every researcher's output near that bound, the optimal allocation and equal division differ little. The value of targeting, the expected output the optimal allocation adds over equal division, vanishes as the budget grows (SI Text 1, Corollary 2), and in every one of several hundred simulated populations it also fell with the budget throughout, as a share of uniform funding's own gain (Fig. S2). This dependence on the budget recurs in every result below.

Proposition~1 concerns one round. Over several rounds a grant also raises the recipient's future capability, and the allocation that maximizes output over the whole horizon need not take the gap form. Numerically it nearly does: over two rounds with complete information, the exact dynamic optimum, free to choose its recipients and how much of the budget to spend in each round, exceeds the gap rule applied round by round with an even split by at most 1.4\% of the research output that funding adds, at every capability tail, compounding rate, and budget we sweep, and its first-round grants correlate with the gap rule's at 0.98 or above (SI Text 1, Table S4). We therefore use the round-by-round gap rule as the complete-information benchmark, and ``optimal allocation'' below refers to it.

\begin{figure}[tp]
\centering
\begin{tikzpicture}
\begin{axis}[
  width=0.82\linewidth, height=6.2cm,
  axis lines*=left,
  xmin=0, xmax=8.6, ymin=0, ymax=7.6,
  xtick=\empty, ytick=\empty,
  x axis line style={-{Stealth[length=4pt]}, line width=0.4pt},
  y axis line style={-{Stealth[length=4pt]}, line width=0.4pt},
  xlabel={resources $R$},
  ylabel={capability $K$},
  ylabel style={font=\small}, xlabel style={font=\small},
  tick label style={font=\small},
  axis line style={line width=0.4pt},
  tick style={line width=0.4pt, black},
  clip=false,
]
% two funding frontiers R = c K, at a smaller budget b_1 (c_1 = 2/3) and a larger budget b_2 (c_2 = 3/2)
\addplot[black, line width=0.5pt, domain=0:7.3, samples=2] ({x*0.6667}, {x});
\addplot[black, line width=0.5pt, domain=0:5.73, samples=2] ({x*1.5}, {x});
\node[font=\small, anchor=west] at (axis cs:4.95,7.25) {frontier at $b_1$};
\node[font=\small, anchor=south east] at (axis cs:8.55,5.85) {frontier at $b_2$};
% grants at the smaller budget: horizontal arrows from (R0, K) to (c_1 K, K)
\addplot[quiver={u=\thisrow{u}, v=\thisrow{v}}, -{Stealth[length=3.5pt]},
         gray, line width=0.5pt] table {fig1-both.dat};
% funded at both budgets (black), funded only at the larger budget (gray), unfunded at both (open)
\addplot[only marks, mark=*, mark size=1.6pt, black] table[x=R, y=K] {fig1-both.dat};
\addplot[only marks, mark=*, mark size=1.6pt, gray!70] table[x=R, y=K] {fig1-b2only.dat};
\addplot[only marks, mark=o, mark size=1.6pt, gray, line width=0.5pt] table[x=R, y=K] {fig1-none.dat};
% direct labels
\node[font=\small, anchor=west] at (axis cs:0.3,6.95) {funded at both budgets};
\node[font=\small, anchor=west, align=left] at (axis cs:3.6,4.6) {funded only at the\\ larger budget};
\node[font=\small, anchor=west, align=left] at (axis cs:4.6,2.4) {unfunded at\\ both budgets};
\end{axis}
\end{tikzpicture}
\caption{The gap rule in one population at two budgets. Each point is a researcher, placed by resources and capability. The two lines are the funding frontiers $R = cK$ at a smaller budget $b_1$ and a larger budget $b_2$; a larger budget raises $c$ and flattens the frontier. Researchers to the left of a frontier are funded at that budget: black points at both budgets, gray points only at the larger budget, open points at neither. Each arrow is a grant at the smaller budget, which moves its researcher exactly to that budget's frontier; at the larger budget every black and gray researcher receives a grant that moves it to the outer frontier (arrows not drawn).}
\label{fig:p1}
\end{figure}


\subsection*{What the funder can learn: track records and peer review}
A real funder observes research output, not capability or resources. Expected output is one number produced by the two jointly, so a researcher of high capability and modest resources and a researcher of modest capability and abundant resources can produce identical records, and the gap rule treats these two most differently: the first may have the largest gap in the population, the second none. Expected output gives only a lower bound on capability, $K_i > \lambda_i/A$, and realized output, a Poisson draw around $\lambda_i$, bounds it less. The problem is worst where the gap is largest. When resources are small relative to capability, expected output is approximately $AR_i$, proportional to resources and nearly independent of capability (Fig.~\ref{fig:p2}\textit{A}), so on small grants output carries little information about capability; only grants comparable to capability itself make records informative about it.

A Bayesian funder relying on records alone therefore allocates far from the gap rule, and further rounds of records improve its allocation only slowly: its expected output falls short of the complete-information optimum by 46\% of what that optimum adds in the first round and still by 19\% in round twenty, while a funder that also holds a single review score for each researcher, observed once at the start, is 8\% short in round one and 3\% by round twenty (Fig.~\ref{fig:p2}\textit{B}). The correlation of the records-only funder's grants with the optimal grants rises from 0.08 to 0.32 over the same rounds, against 0.83 for the review funder in round one. Capabilities compound while resources do not, so each researcher's output approaches the level their resources sustain, and output at that level carries less and less information about capability.

What the funder lacks is a signal of capability separate from resources. Peer review can provide one: a noisy signal, but one whose noise does not come from resources. We model it as each researcher's capability plus Gaussian noise, observed once (Materials and Methods). The value of review is the additional expected output of a funder allocating with the signal over the same funder without it, reported as the share it recovers of the records-only shortfall, the output that complete information would add over records alone. That value depends on the field (Fig.~\ref{fig:p2}\textit{C}). Review is worth most where capability is heavy-tailed and the budget is tight, and under heavy-tailed capability even a signal of moderate informativeness captures most of its value: at the default signal noise, review retains 82\% of its maximum value when the capability distribution's tail parameter is $\alpha_K = 1.3$, 72\% at the default distribution ($\alpha_K = 2$), and 45\% where capability is spread more evenly ($\alpha_K = 3.5$); at three times the default noise, 52\%, 33\%, and 8\%. The three distributions have the same mean capability, so the tail parameter varies how dispersed and how heavy-tailed capability is, not its level. The reason is twofold. Where capability is heavy-tailed, most of the output a funder can add comes from getting resources to the few researchers of unusually high capability, and precisely those researchers are the easiest to identify: they stand far apart from everyone else, so even a noisy signal separates them. Where capability is spread more evenly, no researcher matters much more than another, so review of any informativeness adds little; and where the budget is ample, information about whom to fund has little left to add.

This ordering of review's value by field holds across production functions from Cobb-Douglas to Leontief (SI Text 4), at equal rank correlation between score and capability rather than equal noise, for applicant pools of 25 to 200, for heavier and lighter resource tails, when the review score partly reflects resources, when review is redrawn every round, and across resource signals from nearly exact to nearly uninformative (SI Text 5). Multiplicative review noise, a constant relative error, narrows the heavy-tailed field's advantage without removing it, because a constant relative error is a large absolute error for exactly the researchers who matter most there. What review supplies in every variant is information about capability that output cannot supply on its own; peer review is one source, and the results apply to any other, such as a pilot grant large enough to make output informative. The association between capability and resources changes how much any of this is worth: where the capable already hold the resources, the gaps are small and even, and both targeting and review are worth less. One condition qualifies the value of review: the funder must weight review at its true informativeness. A funder that overtrusts review, treating a noisy signal as reliable, can do worse than one that ignores review altogether, and overtrusting review loses more output than undertrusting it (Fig. S3).

\begin{figure}[tp]
\centering
\begin{minipage}{\linewidth}\noindent\textbf{A}\par\vspace{-0.3em}
\begin{tikzpicture}
\begin{axis}[
  width=0.82\linewidth, height=4.6cm,
  axis lines*=left,
  xmin=0, xmax=6.6, ymin=0, ymax=10.8,
  xtick=\empty, ytick=\empty,
  x axis line style={-{Stealth[length=4pt]}, line width=0.4pt},
  y axis line style={-{Stealth[length=4pt]}, line width=0.4pt},
  xlabel={resources $R$},
  ylabel={expected research output},
  ylabel style={font=\small}, xlabel style={font=\small},
  axis line style={line width=0.4pt},
  clip=false,
]
\addplot[black, line width=0.7pt] table[x=R, y=k30] {fig3-curves.dat};
\addplot[black, line width=0.7pt] table[x=R, y=k10] {fig3-curves.dat};
\addplot[black, line width=0.7pt] table[x=R, y=k3]  {fig3-curves.dat};
\node[font=\small, anchor=west] at (axis cs:6.7,10.0) {$K = 30$};
\node[font=\small, anchor=west] at (axis cs:6.7,7.5)  {$K = 10$};
\node[font=\small, anchor=west] at (axis cs:6.7,4.0)  {$K = 3$};
\end{axis}
\end{tikzpicture}
\end{minipage}\par\vspace{0.9em}
\begin{minipage}{\linewidth}\noindent\textbf{B}\par\vspace{-0.3em}
\begin{tikzpicture}
\begin{axis}[
  width=0.82\linewidth, height=4.6cm,
  axis lines*=left,
  xmin=1, xmax=20, ymin=0, ymax=0.5,
  xtick={1, 5, 10, 15, 20},
  ytick={0, 0.1, 0.2, 0.3, 0.4, 0.5},
  yticklabels={0, 0.1, 0.2, 0.3, 0.4, 0.5},
  xlabel={round},
  ylabel={shortfall (share of complete-information gain)},
  ylabel style={font=\small}, xlabel style={font=\small},
  tick label style={font=\small},
  axis line style={line width=0.4pt},
  tick style={line width=0.4pt, black},
  clip=false,
]
\addplot[black, line width=0.7pt] table[x=round, y=sf4] {fig5-rounds.dat};
\addplot[black, line width=0.7pt] table[x=round, y=sf5] {fig5-rounds.dat};
\node[font=\small, anchor=west] at (axis cs:20.4,0.04) {with review};
\node[font=\small, anchor=west] at (axis cs:20.4,0.20) {no review};
\end{axis}
\end{tikzpicture}
\end{minipage}\par\vspace{0.9em}
\begin{minipage}{\linewidth}\noindent\textbf{C}\par\vspace{-0.3em}
\begin{tikzpicture}
\begin{axis}[
  width=0.82\linewidth, height=4.6cm,
  axis lines*=left,
  xmode=log,
  xmin=0.05, xmax=20, ymin=0, ymax=1,
  xtick={0.05, 0.3, 1, 3, 10, 20},
  xticklabels={0.05, 0.3, 1, 3, 10, 20},
  ytick={0, 0.25, 0.5, 0.75, 1},
  yticklabels={0, 0.25, 0.5, 0.75, 1},
  xlabel={signal noise $\tau_K$},
  ylabel={share of shortfall recovered},
  ylabel style={font=\small}, xlabel style={font=\small},
  tick label style={font=\small},
  axis line style={line width=0.4pt},
  tick style={line width=0.4pt, black},
  clip=false,
]
\addplot[black, line width=0.7pt, error bars/.cd, y dir=both, y explicit, error bar style={line width=0.4pt}, error mark options={rotate=90, mark size=1.2pt, line width=0.4pt}] table[x=tau, y=s13, y error=e13] {fig4-value.dat};
\addplot[black, line width=0.7pt, error bars/.cd, y dir=both, y explicit, error bar style={line width=0.4pt}, error mark options={rotate=90, mark size=1.2pt, line width=0.4pt}] table[x=tau, y=s20, y error=e20] {fig4-value.dat};
\addplot[black, line width=0.7pt, error bars/.cd, y dir=both, y explicit, error bar style={line width=0.4pt}, error mark options={rotate=90, mark size=1.2pt, line width=0.4pt}] table[x=tau, y=s35, y error=e35] {fig4-value.dat};
\node[font=\small, anchor=west] at (axis cs:22,0.13)  {$\alpha_K = 1.3$};
\node[font=\small, anchor=west] at (axis cs:22,0.035) {$\alpha_K = 2$};
\node[font=\small, anchor=west] at (axis cs:22,-0.06) {$\alpha_K = 3.5$};
\end{axis}
\end{tikzpicture}
\end{minipage}\par\vspace{0.9em}
\caption{What the funder can learn. (\textit{A}) On small grants, research output carries little information about capability: expected research output against resources for researchers of capability $K = 3$, $10$, and $30$; the curves coincide where resources are small relative to capability. (\textit{B}) Records do not catch up: the shortfall of each funder's expected research output below the complete-information optimum, as a share of what complete-information funding adds, round by round in the intermediate field; the no-review funder's shortfall falls from 46\% in round one to 19\% by round twenty, while the funder with a single review score observed at the start is at 8\% in round one and 3\% by round twenty. (\textit{C}) Review is worth most where capability is heavy-tailed: the share of the no-review shortfall that the review signal recovers, as review noise $\tau_K$ increases, for three capability distributions with the same mean capability (tail parameters $\alpha_K = 1.3$, $2$, and $3.5$; mean 2 in each); bars are one standard error over 100 simulated populations. Settings: SI Materials and Methods.}
\label{fig:p2}
\end{figure}


\subsection*{Spreading funds: seed grants, lotteries, and concentration}
Two alternatives to targeting recur in funding policy, and each has a rationale. Uniform seed grants give every researcher the same amount, so that a researcher whom review misjudges still receives something; in practice a funder gives out part of its budget as seed grants and targets the rest. Lotteries give grants to applicants chosen at random, usually from among those who pass a review threshold; their advocates argue that review cannot reliably rank the applicants it deems fundable, that review is costly, and that review is biased \citep{FangCasadevall2016, Roumbanis2019}. Both give up part of the value of targeting on purpose, and that value depends on the field: for a funder allocating with records and review, targeting adds more research output than uniform funding's entire gain over no funding in a heavy-tailed field with reliable review and a tight budget, and about a quarter of that gain in the default field at the largest budget we sweep (SI Text 3). What the alternatives save, in review costs and proposal effort, lies outside our model \citep{GrossBergstrom2019}.

Consider seed grants first. The output lost increases faster than the fraction of the budget given out as seed grants: half the budget as seed grants loses more than twice what a quarter loses (Fig.~\ref{fig:p4}\textit{A}). The size of the loss is set by what targeting is worth in the field: a review-informed funder that gives out three quarters of its budget as seed grants forfeits about 11\% of the research output its funding adds in a heavy-tailed field with reliable review (17\% at the smallest budget we sweep), and about 3.5\% in the default field.

Now consider lotteries. The lotteries funders run are partial: review sets a fundable pool, and chance picks the grantees within it. Such a lottery gives up the review ranking within the pool and the certainty of who is funded; we ask what each costs. We evaluate five schemes in a single round (Fig.~\ref{fig:p4}\textit{B}): equal division among the fifth of researchers with the highest review scores (selection by review); equal division among the top two fifths; the partial lottery that some funders run, which funds half of the top two fifths at random with full-size grants; uniform funding; and a lottery over all researchers with full-size grants. In a heavy-tailed field with a tight budget, selection by review gains about three quarters of what the optimal allocation gains when review is reliable and about half when review is very noisy. The partial lottery gains less, and most of the difference comes from ignoring the review ranking within the top two fifths: at the default noise, the partial lottery falls 0.18 short of selection by review, of which the ignored ranking accounts for 0.14 and the randomness for 0.04. The lottery over all researchers gains least, less than uniform funding. In an evenly spread field with a large budget, the ordering reverses: uniform funding gains more than four fifths of what the optimal allocation gains, and every scheme that concentrates grants gains less. A lottery among applicants above a review threshold therefore loses little output only where the review ranking it ignores is worth little: a few hundredths of the optimal allocation's gain in an evenly spread field or when review is very noisy, about a sixth of that gain in a heavy-tailed field with reliable review. In an evenly spread field, the better alternative to review is small grants to many researchers.

The randomness itself costs output for a general reason. A lottery gives each applicant in its pool a chance at a full grant instead of a certain smaller one of the same expected size, and because expected output increases with resources at a diminishing rate, the certain grant produces more. Dividing a pool's money equally therefore produces more expected output than any lottery over the same pool, provided that grants are divisible and that output rises with resources at a diminishing rate at every grant size (SI Text 2, Proposition 2). Where a project needs a minimum viable grant, the first condition fails and a lottery over full awards can beat dividing the same money into grants too small to use. The cost of the randomness is small in our simulations; the cost of the discarded ranking is the one that depends on the field.

Seed grants and lotteries both spread funds more widely than targeting does, and both raise the older question of whether funding should be concentrated among a few researchers or spread across many. The production function does not settle that question: stronger complementarity between capability and resources concentrates the review-informed funder's grants when the budget is tight in a heavy-tailed field, spreads them when the budget is ample, and has a non-monotone effect where capability is evenly spread (Fig. S6; SI Text 3). The production function, the budget, and the field's capability distribution set the answer jointly.

\begin{figure*}[t!]
\centering
\begin{minipage}{\linewidth}\noindent\textbf{A}\par\vspace{-0.3em}
\centering
\begin{tikzpicture}
\begin{axis}[
  width=0.5\linewidth, height=5.2cm,
  axis lines*=left,
  xmin=0, xmax=0.75, ymin=0, ymax=20,
  xtick={0, 0.25, 0.5, 0.75},
  ytick={0, 5, 10, 15, 20},
  xlabel={fraction of the budget given out as seed grants},
  ylabel={research output lost (\% of the gain from funding)},
  ylabel style={font=\small}, xlabel style={font=\small},
  tick label style={font=\small},
  axis line style={line width=0.4pt},
  tick style={line width=0.4pt, black},
  clip=false,
]
\addplot[black, line width=0.7pt, mark=*, mark size=1.2pt] table[x=xseed, y=hb1]  {fig10-floor-cost.dat};
\addplot[black, line width=0.7pt, mark=*, mark size=1.2pt] table[x=xseed, y=hb05] {fig10-floor-cost.dat};
\addplot[black, line width=0.7pt, mark=*, mark size=1.2pt] table[x=xseed, y=hb01] {fig10-floor-cost.dat};
\addplot[black, line width=0.7pt, mark=*, mark size=1.2pt] table[x=xseed, y=db05] {fig10-floor-cost.dat};
\node[font=\small, anchor=west] at (axis cs:0.77,8.11) {heavy-tailed, $b = 2$};
\node[font=\small, anchor=west] at (axis cs:0.77,11.07) {heavy-tailed, $b = 1$};
\node[font=\small, anchor=west] at (axis cs:0.77,17.29) {heavy-tailed, $b = 0.2$};
\node[font=\small, anchor=west] at (axis cs:0.77,3.47) {intermediate, $b = 1$};
\end{axis}
\end{tikzpicture}
\end{minipage}\par\vspace{0.9em}
\begin{minipage}{\linewidth}\noindent\textbf{B}\par\vspace{-0.3em}
\centering
\begin{tikzpicture}
\begin{axis}[
  name=left,
  width=0.32\linewidth, height=4.4cm,
  axis lines*=left,
  xmode=log,
  xmin=0.3, xmax=10, ymin=0, ymax=1,
  xtick={0.3, 1, 3, 10},
  xticklabels={0.3, 1, 3, 10},
  ytick={0, 0.25, 0.5, 0.75, 1},
  yticklabels={0, 0.25, 0.5, 0.75, 1},
  xlabel={review noise $\tau$},
  ylabel={share of the optimal allocation's gain},
  ylabel style={font=\small}, xlabel style={font=\small},
  tick label style={font=\small},
  axis line style={line width=0.4pt},
  tick style={line width=0.4pt, black},
  title={\small heavy-tailed capabilities,\\ tight budget},
  title style={at={(0,1)}, anchor=south west, xshift=-2pt, yshift=-1pt, align=left},
  clip=false,
]
\addplot[black, line width=0.7pt, mark=*, mark size=1.1pt] table[x=tau, y=ss] {fig11-heavy.dat};
\addplot[black, line width=0.7pt, mark=o, mark size=1.3pt] table[x=tau, y=ws] {fig11-heavy.dat};
\addplot[black, line width=0.7pt, mark=triangle*, mark size=1.4pt] table[x=tau, y=sl] {fig11-heavy.dat};
\addplot[black!65, dashed, line width=0.6pt, domain=0.3:10] {0.445};
\addplot[black!65, dotted, line width=0.8pt, domain=0.3:10] {0.378};
\end{axis}
\begin{axis}[
  name=right,
  at={(left.east)}, anchor=west, xshift=30pt,
  width=0.32\linewidth, height=4.4cm,
  axis lines*=left,
  xmode=log,
  xmin=0.3, xmax=10, ymin=0, ymax=1,
  xtick={0.3, 1, 3, 10},
  xticklabels={0.3, 1, 3, 10},
  ytick={0, 0.25, 0.5, 0.75, 1},
  yticklabels={,,,,},
  xlabel={review noise $\tau$},
  xlabel style={font=\small},
  tick label style={font=\small},
  axis line style={line width=0.4pt},
  tick style={line width=0.4pt, black},
  title={\small evenly spread capabilities,\\ ample budget},
  title style={at={(0,1)}, anchor=south west, xshift=-2pt, yshift=-1pt, align=left},
  clip=false,
  legend style={draw=none, fill=none, font=\scriptsize, at={(1.04,0.55)}, anchor=west, legend cell align=left, row sep=1pt},
  legend image post style={mark size=1.1pt},
]
\addplot[black, line width=0.7pt, mark=*, mark size=1.1pt] table[x=tau, y=ss] {fig11-even.dat};
\addlegendentry{top fifth by review}
\addplot[black, line width=0.7pt, mark=o, mark size=1.3pt] table[x=tau, y=ws] {fig11-even.dat};
\addlegendentry{top two fifths by review}
\addplot[black, line width=0.7pt, mark=triangle*, mark size=1.4pt] table[x=tau, y=sl] {fig11-even.dat};
\addlegendentry{partial lottery}
\addplot[black!65, dashed, line width=0.6pt, domain=0.3:10] {0.831};
\addlegendentry{uniform funding}
\addplot[black!65, dotted, line width=0.8pt, domain=0.3:10] {0.414};
\addlegendentry{lottery over all}
\end{axis}
\end{tikzpicture}
\end{minipage}\par\vspace{0.9em}
\caption{Spreading funds. (\textit{A}) The research output lost to uniform seed grants, as a percent of the research output that the funder's grants add over no funding, when a review-informed funder gives out a fraction of its budget as equal grants to all researchers and targets the rest; the loss grows faster than the fraction, and its size is set by what targeting is worth in the field. (\textit{B}) What five funding schemes gain, by field: each scheme's gain over no funding as a share of the optimal allocation's gain, as review noise increases, in a heavy-tailed field with a tight budget (left) and an evenly spread field with a large budget (right). The top fifth and top two fifths by review score receive equal division; the partial lottery funds half of the top two fifths at random with full-size grants. Settings: SI Materials and Methods.}
\label{fig:p4}
\end{figure*}


\section*{Discussion}
Our results isolate one quantity as key to optimal grant funding: the value of targeting researchers with the greatest gap between their capabilities and their resources. Two features of a field determine this value. The first is the funder's budget relative to researchers' own resources: where the budget can relieve most researchers' bottlenecks, the optimal allocation and equal division produce nearly the same research output. The second is how unequally the capability-resource gaps are distributed. Across every setting we ran, the value of targeting is close to a function of the inequality of the optimal grants alone (SI Text 5), and what makes the gaps unequal is mainly how unequal capability is; the association between capability and resources works through the same channel, since where the capable already hold the resources the gaps are small and even. Two results lie outside this quantity. How widely a review-informed funder should spread its grants depends, in addition, on how strongly capability and resources complement each other (Fig. S6); and the timing of spending over the horizon, analyzed in SI Text 6, adds little research output beside the choice of recipients in every setting we test, within the limits stated there.

Within the model, no funding mechanism is best in every field, so a funder choosing among peer review, seed grants, and lotteries should begin by locating its field on the two features (Fig.~\ref{fig:regime}). The first is observable: a funder's annual awards to a field against the field's research spending from every other source. The second is not, because the gaps depend on capabilities the funder cannot see and on the budget itself. Three observable quantities stand in for it. The inequality of research output in the field is the natural proxy for the inequality of capability, though where grants are small output tracks resources rather than capability, so this proxy leans toward the resource distribution. The dispersion of review scores over the whole applicant pool is a second, with the same caveat about noise. The most direct is the response of output to grants large enough to relieve the bottleneck (Fig.~\ref{fig:p2}\textit{A}): a program of substantial pilot grants measures the gaps, whereas small seed grants do not. How existing funding is distributed completes the estimate: where resources already sit with the capable, the gaps are smaller and more even than capability inequality alone would imply.

Fig.~\ref{fig:regime} places four kinds of funding program on the budget axis by order of magnitude, from 2023 and 2024 budget figures (SI Materials and Methods). A national biomedical agency that supplies a large share of its field's research support sits near a budget equal to the field's other resources; a national science agency that funds its fields more thinly near a tenth of that; a continental excellence council near a few hundredths; and a private foundation funding one field among many at a hundredth or below. Where a field lies on the inequality axis is the estimate each funder must make, so the placements say which regime a funder's budget puts it in, not where it sits within the regime. The regimes differ. A foundation or council with a small budget in a heavy-tailed field faces the largest value of targeting per dollar: review directed at identifying the few researchers of unusually high capability is worth most, even noisy review recovers most of that value, seed grants give up the most output, and a lottery among applicants above a review threshold gives up about a sixth of what selection by review gains. An agency whose budget approaches the field's other resources, in a field where capability is evenly spread, faces a small value of targeting: review of any informativeness adds little, seed grants and lotteries give up little, and the model's best use of the budget is small grants to many researchers.

% PNAS Fig. 4 (Discussion): the regime map, two panels; copy of fig16-regime.tex with SI refs. x = budget relative to the field's
% resources (b, log scale); y = capability inequality (Gini coefficient of the Pareto
% capability distribution, 1/(2 alpha_K - 1), with alpha_K in parentheses). (A) value of
% targeting: the complete-information optimum's gain over uniform funding as a share of
% uniform funding's gain. (B) value of review at the default noise, as a share of the research
% output that review-informed funding adds. Bars above the panels: four kinds of program placed on the
% budget axis by order of magnitude (SI Materials and Methods). Data: fig16-targ.dat, fig16-rev.dat from
% regime_m.R (M = 1000) via make_regime.py; contour label positions computed from the data.
\begin{figure*}[t]
\centering
\pgfplotsset{keep level/.style={x filter/.append code={\edef\tempa{\thisrow{level}}\edef\tempb{#1}\ifx\tempa\tempb\else\def\pgfmathresult{inf}\fi}}}
\newcommand{\programmarks}{%
\draw[line width=0.5pt] (axis cs:0.01,0.745) -- (axis cs:0.03,0.745);
\draw[line width=0.5pt] (axis cs:0.04,0.745) -- (axis cs:0.08,0.745);
\draw[line width=0.5pt] (axis cs:0.2,0.745) -- (axis cs:0.45,0.745);
\draw[line width=0.5pt] (axis cs:1.5,0.745) -- (axis cs:3,0.745);
\node[font=\scriptsize, anchor=south] at (axis cs:0.0173,0.755) {foundation};
\node[font=\scriptsize, anchor=south] at (axis cs:0.0566,0.81) {excellence council};
\node[font=\scriptsize, anchor=south] at (axis cs:0.3,0.755) {science agency};
\node[font=\scriptsize, anchor=south] at (axis cs:2.12,0.81) {biomedical agency};}
\begin{tikzpicture}
\begin{axis}[
  name=panelA,
  scale only axis, width=4.1cm, height=5.4cm,
  axis lines*=left,
  xmode=log,
  xmin=0.01, xmax=3, ymin=0.1, ymax=0.72,
  xtick={0.01, 0.1, 1},
  xticklabels={0.01, 0.1, 1},
  xminorticks=false,
  ytick={0.167, 0.333, 0.625},
  yticklabels={evenly spread, intermediate, heavy-tailed},
  xlabel={budget relative to the field's resources, $b$},
  ylabel={capability inequality},
  ylabel style={font=\small}, xlabel style={font=\small, at={(rel axis cs:1.14,-0.1)}, anchor=north},
  tick label style={font=\small},
  axis line style={line width=0.4pt},
  tick style={line width=0.4pt, black},
  clip=false,
]
\programmarks
\foreach \lev in {0.1, 0.25, 0.5, 1, 2} {
  \addplot[black, line width=0.6pt, empty line=jump, keep level=\lev] table[x expr=10^\thisrow{x}, y=y] {fig16-targ.dat};
}
\node[font=\scriptsize, anchor=west] at (axis cs:3,0.178) {0.1};
\node[font=\scriptsize, anchor=west] at (axis cs:3,0.304) {0.25};
\node[font=\scriptsize, anchor=west] at (axis cs:3,0.446) {0.5};
\node[font=\scriptsize, anchor=west] at (axis cs:3,0.639) {1};
\node[font=\scriptsize, anchor=west] at (axis cs:0.72,0.695) {2};
\node[font=\bfseries, anchor=south west] at (rel axis cs:-0.36,1.2) {A};
\end{axis}
\begin{axis}[
  at={(panelA.east)}, anchor=west, xshift=1.25cm,
  scale only axis, width=4.1cm, height=5.4cm,
  axis lines*=left,
  xmode=log,
  xmin=0.01, xmax=3, ymin=0.1, ymax=0.72,
  xtick={0.01, 0.1, 1},
  xticklabels={0.01, 0.1, 1},
  xminorticks=false,
  ytick={0.167, 0.333, 0.625},
  yticklabels={},
  xlabel={}, ylabel={},
  tick label style={font=\small},
  axis line style={line width=0.4pt},
  tick style={line width=0.4pt, black},
  clip=false,
]
\programmarks
\foreach \lev in {0.05, 0.1, 0.2, 0.3} {
  \addplot[black, line width=0.6pt, empty line=jump, keep level=\lev] table[x expr=10^\thisrow{x}, y=y] {fig16-rev.dat};
}
\node[font=\scriptsize, anchor=west] at (axis cs:3,0.199) {0.05};
\node[font=\scriptsize, anchor=west] at (axis cs:3,0.277) {0.1};
\node[font=\scriptsize, anchor=west] at (axis cs:3,0.435) {0.2};
\node[font=\scriptsize, anchor=west] at (axis cs:3,0.693) {0.3};
\node[font=\bfseries, anchor=south west] at (rel axis cs:-0.1,1.2) {B};
\end{axis}
\end{tikzpicture}
\caption{Where a funder stands. (\textit{A}) The value of targeting: the complete-information optimum's gain over uniform funding, as a share of uniform funding's gain over no funding. (\textit{B}) The value of review at the default noise, as a share of the research output that review-informed funding adds. Both against the funder's budget relative to the field's resources (horizontal) and the inequality of capability (vertical: the Gini coefficient of capability, with ticks at the paper's three fields, $\alpha_K = 1.3$, $2$, and $3.5$). Contours are labeled by value. The bars above each panel place four kinds of program on the budget axis by order of magnitude, from 2023 and 2024 annual budgets (SI Materials and Methods); where a program's field lies on the vertical axis is the estimate each funder must make. Two rounds, $\tau_K = 1$, 50 populations per cell.}
\label{fig:regime}
\end{figure*}


Three further implications concern peer review. First, most of review's value comes from separating the few researchers of unusually high capability from the rest, so a review system that spends most of its effort ranking the many applications of comparable merit spends that effort where the model locates the least output. Second, the finding that percentile scores barely predict productivity comes from funded NIH grants with scores in the top fifth \citep{Fang2016}, and the studies that find modest predictive power also use funded grants \citep{LiAgha2015}; both measure prediction among applicants review has already selected, where the model expects review to add least, so a weak association there is consistent with review of considerable value over the whole applicant pool. Third, a funder that treats review as more reliable than it is loses more output than a funder that treats review as noisier than it is, so a funder uncertain about review's accuracy should err toward undertrust.

The model holds several things fixed, and its results hold within those limits. There is one funder, and researchers do not respond strategically, so competition among funders, proposal effort, risk choice, and the cost of review are outside the model; the contest models of grant competition cover that ground \citep{GrossBergstrom2019, GrossBergstrom2025}, and the mechanism-design treatment of grants by Carnehl et al.~\cite{Carnehl2024}, in which the funder must elicit information from applicants who bear the cost of applying, complements a model in which the funder's information arrives for free. Sorzano and Pueche-Granados~\cite{Sorzano2026} also treat funding as a decision under heavy-tailed returns and find more room for lotteries than we do; their heavy tail is in observed impact, whereas ours decomposes output into latent capability and resources, and that decomposition is what gives review its value here. Output is a single number per round, so projects do not differ in kind and exploring new directions has no value in the model \citep{AvinBJPS2019}; nor do a researcher's proposals differ in quality, so this is a model of investigator-level funding. Grants are divisible (Results). The review signal is noisy but unbiased; bias and conservatism in review, which are among the strongest arguments for lotteries, are not modeled. Capability is a state, not a trait: it holds whatever a researcher has accumulated that raises output, and doing research adds to it. Capability persists and grows while resources recur and do not persist, so a grant's lasting effect enters only through capability; durable equipment, data, and infrastructure violate this asymmetry, and the results about compounding and timing (SI Text 6) depend on it while the one-round results do not. The compounding here is of capability through doing research, and differs from the cumulative advantage in which winning a grant raises the chance of winning the next \citep{BolVaanRijt2018}. None of the numbers we report is calibrated to a real funder; what applies to a real field is the qualitative features of the dynamics that organize the results.

The model makes four predictions that data on funded research can test. The effect of winning a grant on output should be larger for recipients whose other resources are small relative to their capability, so a randomized trial of funding, such as the one that found no clear average effect \citep{Barnett2024}, should find its effect concentrated among under-resourced recipients; a null average is what the model predicts where most recipients are already well resourced. Measures that stand in for capability, such as prior record or review score, should predict research output better the larger a grant is relative to the recipient's other resources. How well review scores predict output, measured over the whole applicant pool rather than among funded grants, should be higher in fields whose output distributions have heavier tails. And where funders run lotteries beside review, the output lost to the lottery should be larger in heavier-tailed fields.

A funder with a fixed budget and imperfect knowledge of researchers faces one problem, and our model gives it one answer. The output-maximizing allocation funds the gap between each researcher's capability and their resources. Records alone underdetermine that gap, because output does not identify which input binds; a signal of capability separate from resources does, and peer review is one source of it. What that signal is worth, and how seed grants and lotteries perform, turn on two features of a field, and the map of those features has regimes in which the answers differ. A funder deciding among review, lotteries, and seed grants should begin by asking which kind of field it is funding.


\matmethods{\subsection*{Model} A funder faces $n$ researchers over $T$ funding rounds. Researcher $i$ has capability $K_i$ and resources $R_i$; initial values are drawn from Pareto distributions with tail parameters $\alpha_K$ and $\alpha_R$ (smaller values give more unequal fields; below 2 the variance is infinite), coupled through a Gaussian copula with parameter $\rho$ ($\rho = 0.5$ and $0.8$ give rank correlations of 0.48 and 0.79; a Pearson correlation is undefined when the variances are infinite). In each round, researcher $i$ produces research output drawn from a Poisson distribution with mean $\lambda_i = AK_iR_i/(K_i + R_i)$, proportional to the harmonic mean of the two inputs, where the productivity constant $A$ sets the units of output (the simulations set $A = 1$). Between rounds capability compounds, $K_i \leftarrow K_i + \epsilon\,\lambda_i/A$, at rate $\epsilon$. Resources do not accumulate: $R_i = R_{i0} + g_i$, where the baseline $R_{i0}$ recurs every round and the grant $g_i$ is consumed in the round it is given. The unit of analysis is the researcher, not a researcher-project pair. The funder's total budget over the horizon is $b\,n\,\mathbb{E}[R_0]$, so at budget scale $b = 1$ it equals the community's expected baseline resources for one round; the budget does not depend on the horizon.

\subsection*{Information and strategies} The funder knows the distributions but observes no researcher's capability or resources. It observes each researcher's realized output as it accumulates; a resource signal $R_{i0} + \eta_i$, observed once, with $\eta_i$ Gaussian with standard deviation $\tau_R$; and a review signal $K_i + \xi_i$, observed once at the start by the strategies that use it, with $\xi_i$ Gaussian with standard deviation $\tau_K$ and independent across researchers, so a reviewer's misjudgment persists (SI Text 5 redraws review every round). Bayesian strategies hold a posterior over each researcher's capability and baseline resources, computed by importance sampling with $M$ particles drawn from the prior ($M = 1{,}000$ for Fig.~\ref{fig:p2}\textit{B} and \textit{C}, Fig.~\ref{fig:regime}, and the sensitivity analysis, and $200$ elsewhere; convergence in $M$ is checked in SI Materials and Methods), update it after each round, and allocate each dollar to the researcher whose expected output, given the posterior, rises most with it. The myopic funder maximizes the current round's expected output; the forward-looking funder plans the remaining budget over the remaining horizon, treating future output as equal to its expectation, and re-plans each round (SI Text 6). The complete-information benchmark applies the gap rule to each round's share of the budget with the true capabilities and resources. Three baselines complete the set: no funding, against which every result is normalized; uniform funding; and track-record funding, which allocates in proportion to the previous round's output (SI Materials and Methods, Table S1).

\subsection*{Parameters and scoring} Tail parameters range from 1.3 to 3.5, a range that includes the Pareto tails near 1 to 2 estimated for research productivity \citep{Lotka1926, Nicholls1989} and extends well beyond them, since productivity is an imperfect guide to capability (Results); the budget scale $b$ from 0.01 to 3; the compounding rate $\epsilon$ from near zero to 0.85; review noise $\tau_K$ from 0.05 to 20; the horizon from one round to ten. Default values are $T = 2$, $n = 50$, $\epsilon = 0.1$, $b = 1$, $\alpha_K = \alpha_R = 2$, $\tau_K = \tau_R = 1$, $\rho = 0$. Each strategy is scored by the sum of expected outputs along its run, and strategies are compared on identical simulated populations with identical signal and output draws, so comparisons are paired. Results are reported as shares of a named comparison, most often the research output that funding adds: a strategy's expected output minus the expected output under no funding. Proofs are in SI Text 1 and SI Text 2; the settings behind each figure are in SI Materials and Methods. Code and data are available in the project repository. }
\showmatmethods{}
\acknow{This work was supported by the John Templeton Foundation [grant number TODO].}
\showacknow{}
\bibliography{references}
\end{document}
